% =====================================================================
%  FORMULACION UNIFICADA
%  Salida coronaria no-Newtoniana con ley de tubo colapsable
%  Modelo acoplado LV 0D <-> coronaria 3D (CoupledLV0DCoronary3D)
%  Requiere: amsmath, amssymb, booktabs
% =====================================================================
\documentclass[11pt]{article}
\usepackage{amsmath,amssymb}
\usepackage{booktabs}
\usepackage[margin=2.3cm]{geometry}
\newcommand{\dd}[2]{\frac{\partial #1}{\partial #2}}
\newcommand{\ptm}{p_{\mathrm{tm}}}
\title{Salida coronaria no-Newtoniana con ley de tubo colapsable:\\
       formulaci\'on unificada}
\date{}
\begin{document}\maketitle

\begin{abstract}
\noindent
Se presenta una reformulaci\'on de la condici\'on de salida coronaria del
modelo acoplado 0D--3D. La versi\'on anterior conten\'ia la bomba
intramioc\'ardica pero no la modulaci\'on de conductancia; la revisi\'on
recibida propon\'ia la modulaci\'on de conductancia pero eliminaba la bomba.
Ambas son medios modelos. Aqu\'i los dos mecanismos se reducen a un \'unico
enunciado constitutivo: el radio de los vasos intramioc\'ardicos obedece una
ley de tubo colapsable $r=r(\ptm)$, y ese mismo radio alimenta la ley de flujo
(modulaci\'on de conductancia) y el volumen anat\'omico $N\pi r^{2}L$ (bomba).
De ah\'i se siguen, sin t\'erminos adicionales, la restricci\'on $V\ge0$, el
volumen no tensionado, la compliancia y la constante de vaciado. La
calibraci\'on se reordena en torno a \emph{tiempos de tr\'ansito}, con lo que
desaparece el residuo ajustado $L_a$ y los calibres pasan a ser predicciones
verificables. El n\'umero de constantes libres de la condici\'on de salida baja
de 15 a 12 y desaparecen las tres menos defendibles ($p_0$, $\alpha_a$,
$\tau_f$). En una prueba 0D independiente, el cociente sist\'ole/di\'astole del
flujo de entrada pasa de $1.97$ a $0.54$, dentro del rango fisiol\'ogico
$0.4$--$0.6$; \'este era el defecto cualitativo pendiente.
\end{abstract}

\section{Qu\'e estaba mal, y d\'onde}
\label{sec:diag}

Conviene fijar el diagn\'ostico antes de la formulaci\'on, porque determina
qu\'e se conserva y qu\'e se sustituye.

\paragraph{El patr\'on f\'asico estaba invertido.} En la simulaci\'on
disponible, el flujo medio de entrada durante la s\'istole era $2.22$ veces el
diast\'olico, con el m\'aximo de entrada arterial en fase $0.226$ del ciclo,
pr\'acticamente simult\'aneo al m\'aximo de salida venosa (fase $0.232$). La
circulaci\'on coronaria izquierda es diast\'olica: el cociente
sist\'ole/di\'astole del flujo de entrada vale $0.4$--$0.6$. Un modelo sin
modulaci\'on de conductancia no puede reproducirlo, porque el \'unico
mecanismo sist\'olico disponible, la contrapresi\'on $p_c$ generada por la
bomba, queda compensado por el ascenso simult\'aneo de $p_{ar}$: ambas
presiones suben juntas y la ca\'ida motriz $p_{ca}-p_c$ apenas se reduce. Lo
que falta no es amplitud de bomba sino resistencia dependiente de la
compresi\'on.

\paragraph{La reserva expulsable era nula.} La presi\'on transmural registrada
recorr\'ia el intervalo $[-8560,\,1146]$~Pa. Con la ley \emph{softplus}
$V=Cp_0\ln(1+e^{p/p_0})$ eso significa que el lecho se vac\'ia por completo en
cada latido: el compartimento se queda literalmente sin sangre a los pocos
milisegundos de la s\'istole y permanece as\'i. El error no est\'a en $C$
---$C_{tot}=3.2\times10^{-10}$~m$^3$/Pa es compatible con las
identificaciones publicadas--- sino en que la ley \eqref{eq:vcol-old} asigna
volumen \emph{cero} a presi\'on transmural cero,
\begin{equation}
  V(p)=C\,p_0\ln\bigl(1+e^{p/p_0}\bigr)\;\longrightarrow\;0
  \quad\text{cuando}\quad \ptm\to0 ,
  \label{eq:vcol-old}
\end{equation}
de modo que todo el dep\'osito vale $C\times1740~\text{Pa}=0.56$~mL, mientras
que la propia anatom\'ia del modelo implica bastante m\'as: el \'arbol venular
solo ya contiene $N_v\pi r_v^2L_v=QL_v/\bar v_v\approx10$~mL. Un vaso real a
presi\'on transmural nula conserva su volumen no tensionado y s\'olo se vac\'ia
cuando la pared pandea a $\ptm$ apreciablemente negativa. La correcci\'on no es
retocar $p_0$ ni $C$: es sustituir una sigmoide por el origen por una ley de
tubo colapsable con $V_0$ anat\'omico.

\paragraph{La calibraci\'on era circular.} $L_a$ era el residuo del presupuesto
de presi\'on ($\approx46$~mm, un ``camino efectivo'' sin correlato
morfom\'etrico), $C_{tot}$ proven\'ia de una identificaci\'on externa,
$L_v=1.5$~cm se impon\'ia sobre una v\'enula de $30~\mu$m, y $C_a/C=0.10$
resultaba ser unas siete veces el cociente de esa misma identificaci\'on
canina, se\~nal de que $C_a$ se hab\'ia dimensionado para satisfacer la
condici\'on de estabilidad y no por medida. Adem\'as $Z_c$ y $C_a$, que son la
misma elasticidad vista de dos maneras, se fijaban por separado y de forma
mutuamente inconsistente.

\paragraph{Lo art\'ificial estaba en la presi\'on tisular.} Los tres
par\'ametros sin medida independiente eran $\alpha$, $\alpha_a=0.01$ y
$\tau_f=0.02$~s. El filtro exist\'ia \'unicamente para acotar $\dot p_{im}$
frente a una fuente expl\'icita $Q_{im}=C\dot p_{im}$; una vez que esa fuente
desaparece, el filtro no tiene funci\'on.

\paragraph{Atribuci\'on.} Por completitud: Spaan, Breuls \& Laird (1981) es el
trabajo de la \emph{bomba} intramioc\'ardica, y Spaan (1985) lo extiende a la
presi\'on de flujo cero mediante la compliancia intramioc\'ardica; la
modulaci\'on de conductancia y la contrapresi\'on pertenecen a la l\'inea de la
\emph{cascada vascular}, Downey \& Kirk (1975) y Westerhof, Sipkema \& van Huis
(1983). La posici\'on moderna (Westerhof et al., \emph{Physiol Rev} 2006;
Algranati, Kassab \& Lanir 2010) es que ambos mecanismos operan y que ninguno
por s\'i solo reproduce la onda f\'asica.

\section{Principio rector}

\begin{quote}
\emph{La geometr\'ia representa anatom\'ia; la resistencia hidr\'aulica
representa la red vascular no resuelta; y la pared del vaso, una sola ley
constitutiva, produce a la vez la resistencia y el almacenamiento.}
\end{quote}

\noindent
La tercera cl\'ausula es la novedad. Los dos mecanismos discutidos no son dos
t\'erminos que haya que sumar ni dos opciones entre las que haya que elegir:
son dos consecuencias de que el calibre dependa de la presi\'on transmural. Por
eso esta formulaci\'on tiene \emph{menos} constantes que cualquiera de los dos
mecanismos modelado por separado.

\section{Reolog\'ia y ley de tubo}

\begin{equation}
  \mu(\dot\gamma)=\mu_\infty+(\mu_0-\mu_\infty)
  \Bigl[1+(\lambda\dot\gamma)^{a}\Bigr]^{\frac{n-1}{a}} ,
  \label{eq:cy}
\end{equation}
la misma en el dominio 3D y en los elementos 0D, de modo que el acoplamiento es
constitutivamente consistente. Para un tubo recto de radio $r$ y longitud $L$,
con $\tau_w=r|\Delta p|/2L$ y $\mu(\dot\gamma_w)\dot\gamma_w=\tau_w$,
\begin{equation}
  q(\Delta p;L,r)=\frac{\pi r^3}{\tau_w^{3}}
  \int_{0}^{\dot\gamma_w}\!\dot\gamma^{3}\,\mu^2\,
  \bigl(\mu+\dot\gamma\,\mu'\bigr)\,d\dot\gamma .
  \label{eq:wrms}
\end{equation}

\paragraph{Las dos derivadas exactas.} La primera ya estaba:
\begin{equation}
  \dd{q}{\Delta p}=\frac{\pi r^3\,\dot\gamma_w-3\,q}{\Delta p} .
  \label{eq:dqdp}
\end{equation}
La segunda es nueva y hace viable todo lo dem\'as. Escribiendo
\eqref{eq:wrms} en funci\'on de $\tau_w$ se observa que el prefactor
\emph{no contiene el radio},
\[
  \frac{\pi r^{3}}{\tau_w^{3}}=\frac{8\pi L^{3}}{\Delta p^{3}},
  \qquad\text{luego}\qquad
  q=\frac{8\pi L^{3}}{\Delta p^{3}}\,I(\dot\gamma_w),
\]
y toda la dependencia en $r$ pasa por $\dot\gamma_w$. Derivando
$I'(\dot\gamma_w)=\dot\gamma_w^{3}\mu_w^{2}(\mu_w+\dot\gamma_w\mu_w')$ y usando
$(\mu+\dot\gamma\mu')\,\partial\dot\gamma_w/\partial r=\tau_w/r$ se obtiene, tras
cancelar,
\begin{equation}
  \boxed{\;\dd{q}{r}=\pi r^{2}\,\dot\gamma_w\;}
  \label{eq:dqdr}
\end{equation}
exacta para cualquier $\mu(\dot\gamma)$ generalizada newtoniana, y reducible a
$4q/r$ en el l\'imite de Poiseuille. Su importancia pr\'actica es que
$\dot\gamma_w$ y $q$ ya se han calculado: convertir el radio en variable de
estado deja el Jacobiano \emph{exacto} y el coste por iteraci\'on de Newton
\emph{sin cambios}.

\paragraph{Relaci\'on de diagn\'ostico.} Con $Q=\pi r^2\bar v$,
\begin{equation}
  \dot\gamma_w=\frac{4Q}{\pi r^3}=\frac{4\bar v}{r},
  \label{eq:shear}
\end{equation}
independiente de $L$: el radio y la velocidad fijan el r\'egimen reol\'ogico
mientras que la longitud s\'olo escala la resistencia. Esta identidad gobierna
el dise\~no de la calibraci\'on.

\section{La ley de tubo colapsable}
\label{sec:tubo}

Cada \'arbol intramioc\'ardico es un lecho de $N$ vasos id\'enticos de
\'area no tensionada $A_u$ y longitud de camino $L$. El \'area obedece la ley
cl\'asica de tubo colapsable (Shapiro; Kamm \& Shapiro)
\begin{equation}
  \ptm = K\Bigl[\,a^{m}-a^{-n}\,\Bigr],
  \qquad a=\frac{A}{A_u},
  \qquad m=10,\quad n=\tfrac32 ,
  \label{eq:tubo}
\end{equation}
con $K=\dfrac{E}{12(1-\nu^{2})}\Bigl(\dfrac{h}{r}\Bigr)^{3}$ la rigidez de
flexi\'on de la pared. La rama $a^{m}$ es la pared en tracci\'on, que se
endurece muy r\'apido; la rama $-a^{-n}$ es el pandeo en compresi\'on. Se
invierte para $a$ con un Newton salvaguardado por bisecci\'on sobre $s=\ln a$,
donde el residuo es mon\'otono, y la derivada es exacta por derivaci\'on
impl\'icita:
\begin{equation}
  \frac{da}{d\ptm}=\frac{1}{K\bigl(m\,a^{m-1}+n\,a^{-n-1}\bigr)} .
  \label{eq:dadp}
\end{equation}

\paragraph{Lo que la ley concede gratis.} Tres propiedades que antes se
imponan con t\'erminos separados:
\begin{itemize}
\item $a>0$ para todo $\ptm$ finito, luego $V\ge0$ \emph{por construcci\'on},
      sin penalizaci\'on ni saturaci\'on ad hoc;
\item $a=1$ exactamente en $\ptm=0$: el vaso conserva su volumen no
      tensionado a presi\'on transmural nula y s\'olo se vac\'ia al pandear.
      \'Este es el defecto de \eqref{eq:vcol-old} que se corrige;
\item la compliancia $dV/d\ptm$ deja de ser un dato y pasa a ser
      $N L A_u\,da/d\ptm$, es decir, la rigidez de pared actuando sobre el
      volumen anat\'omico. Se reporta como predicci\'on.
\end{itemize}

\paragraph{Un solo punto de evaluaci\'on.} La resistencia de un \'arbol
coronario se concentra en sus generaciones terminales, que son intramioc\'ardicas
y est\'an en el extremo aguas abajo del \'arbol arteriolar y aguas arriba del
venular; ambas, por tanto, en el nodo $p_c$. Se evaluan los dos \'arboles en la
misma presi\'on transmural
\begin{equation}
  \ptm = p_c-p_{im} .
  \label{eq:ptm}
\end{equation}
Esto no es una simplificaci\'on de conveniencia: es la misma hip\'otesis que ya
sosten\'ia el t\'ermino $V(p_c-p_{im})$ de la formulaci\'on anterior, extendida
ahora tambi\'en a los radios. La consecuencia es que \emph{el balance de nodo no
cambia}; s\'olo cambia de qu\'e depende $V$ y de qu\'e dependen los radios.

\section{Presi\'on tisular}

\begin{equation}
  p_{im}=\alpha\,p_{LV} .
  \label{eq:pim}
\end{equation}
Nada m\'as. El t\'ermino de acortamiento $\alpha_a\tau_c$ y el retardo de primer
orden $\tau_f$ se eliminan: ninguno ten\'ia medida independiente, y el retardo
exist\'ia solamente para acotar $\dot p_{im}$ frente a la fuente expl\'icita
$Q_{im}=C\dot p_{im}$, que ya no existe. Ahora $p_{im}$ entra \'unicamente a
trav\'es de \eqref{eq:tubo}, resuelta impl\'icitamente, y el volumen que puede
desplazar est\'a acotado por el contenido anat\'omico del lecho; una $p_{LV}$
casi discontinua la absorbe la f\'isica, no un filtro. La verificaci\'on 0D lo
confirma: reintroducir el filtro con $\tau_f$ hasta $0.04$~s modifica el
cociente pico/media venoso de $4.27$ a $3.73$ y empeora el cociente
sist\'ole/di\'astole, de modo que no explica nada que la f\'isica no explique
mejor.

$\alpha$ es una transmisi\'on promediada sobre el espesor de pared: la presi\'on
tisular va de $\approx p_{LV}$ en el subendocardio a una fracci\'on peque\~na en
el subepicardio, y el promedio ponderado por volumen es del orden de
$0.5$--$0.7$. Es la \'unica constante a la que el modelo es genuinamente
sensible (\S\ref{sec:sens}), y as\'i debe declararse.

\section{Balance de nodo y discretizaci\'on}

\[
\text{salida 3D}\;\rightarrow\;
\underbrace{\text{conducto}}_{r_p\ \text{de la malla}}\;\rightarrow\;
\underbrace{p_{ca},\;C_a}_{\text{epic\'ardico}}\;\rightarrow\;
\underbrace{\text{\'arbol arteriolar}}_{N_a,\,L_a,\,r_a(\ptm)}\;\rightarrow\;
\underbrace{p_c\;[\,p_{im}(t)\,]}_{\text{microvascular}}\;\rightarrow\;
\underbrace{\text{\'arbol venular}}_{N_v,\,L_v,\,r_v(\ptm)}\;\rightarrow\; P_{RA}
\]

\noindent
Volumen almacenado y balance:
\begin{align}
  V(\ptm) &= N_a L_a A_{u,a}\,a_a(\ptm)+N_v L_v A_{u,v}\,a_v(\ptm),
  \label{eq:V}\\[2pt]
  C_a\,\frac{dp_{ca}}{dt} &= Q-q_a,
  \qquad
  \frac{d}{dt}V\bigl(p_c-p_{im}\bigr)=q_a-q_v .
  \label{eq:balance}
\end{align}
El balance \eqref{eq:balance} es, t\'ermino a t\'ermino, el de la formulaci\'on
anterior. Conviene decirlo expl\'icitamente porque es donde la revisi\'on
recibida se equivocaba: $\tfrac{d}{dt}V(p_c-p_{im})=q_a-q_v$ \emph{es} la bomba
intramioc\'ardica ---desarrollando,
$C_{\text{eff}}\dot p_c=q_a-q_v+C_{\text{eff}}\dot p_{im}$, y el \'ultimo
t\'ermino es exactamente la antigua fuente expl\'icita---, y no es un
a\~nadido artificial sino conservaci\'on de volumen para un compartimento cuyo
entorno est\'a presurizado. Lo \'unico que se a\~nade es una dependencia: los
radios, y por tanto $V$, dependen de $\ptm$.

Con Euler impl\'icito, inc\'ognitas $(p_{ca},p_c)$ y Newton $2\times2$:
\begin{align}
  R_1&=\frac{C_a}{\Delta t}\bigl(p_{ca}-p_{ca}^{n}\bigr)+q_a-Q=0,
  &q_a&=N_a\,q\bigl(p_{ca}-p_c;\,L_a,\,r_a(\ptm)\bigr),
  \label{eq:R1}\\
  R_2&=\frac{V(\ptm)-V^{n}}{\Delta t}-q_a+q_v=0,
  &q_v&=N_v\,q\bigl(p_c-P_{RA};\,L_v,\,r_v(\ptm)\bigr) .
  \label{eq:R2}
\end{align}
Con $r=\sqrt{A_u a/\pi}$ y por tanto
$dr/d\ptm=A_u(da/d\ptm)/(2\pi r)$, y usando \eqref{eq:dqdp} y \eqref{eq:dqdr},
\begin{equation}
  J=\begin{pmatrix}
    \dfrac{C_a}{\Delta t}+N_a q_a' & \ \ \partial_{p_c}q_a\\[8pt]
    -N_a q_a' & \ \ \dfrac{C_{\text{eff}}}{\Delta t}-\partial_{p_c}q_a
      +\partial_{p_c}q_v
  \end{pmatrix},
  \label{eq:jac}
\end{equation}
con
\begin{equation}
  \partial_{p_c}q_a=N_a\Bigl[-q_a'+\pi r_a^{2}\dot\gamma_{w,a}\,\frac{dr_a}{d\ptm}\Bigr],
  \quad
  \partial_{p_c}q_v=N_v\Bigl[\ \ q_v'+\pi r_v^{2}\dot\gamma_{w,v}\,\frac{dr_v}{d\ptm}\Bigr],
  \quad
  C_{\text{eff}}=\dd{V}{\ptm} .
\end{equation}
El Jacobiano deja de ser sim\'etrico, lo que es irrelevante en un sistema
$2\times2$ resuelto directamente, y todas sus entradas son exactas: se ha
comprobado frente a diferencias finitas con error relativo $4\times10^{-7}$. Se
a\~nade una b\'usqueda lineal por retroceso sobre $|R_1|+|R_2|$, porque la rama
en tracci\'on de \eqref{eq:tubo} es muy r\'igida y un paso completo puede
sobrepasar hacia la rama colapsada; con ella el solver converge en $3.6$
iteraciones de media y $6$ como m\'aximo.

\paragraph{Impedancia caracter\'istica y compliancia epic\'ardica.} Son la misma
elasticidad vista dos veces (Moens--Korteweg, $c=1/\sqrt{\rho D}$), y en la
formulaci\'on anterior se fijaban por separado. Aqu\'i comparten la \'unica
constante $c$:
\begin{equation}
  Z_c=\frac{\rho\,c}{A_p},
  \qquad
  C_a=\frac{A_p L_p}{\rho\,c^{2}},
  \qquad\text{de donde}\qquad
  Z_c\,C_a=\frac{L_p}{c},
  \label{eq:zc}
\end{equation}
el tiempo de tr\'ansito de onda del conducto. La presi\'on de salida sigue
siendo $p_{out}=p_{ca}+\Delta P_p(Q)+Z_cQ$.

\paragraph{Estabilidad del acoplamiento 3D--0D.} La presi\'on de salida se
aplica con un retraso de un paso, luego el 3D debe ver un nodo capacitivo y no
la resistencia arteriolar dominante: $\Delta t/C_a<Z_{3D}\approx\rho L_{3D}/(A\Delta t)$.
Sustituyendo \eqref{eq:zc}, el \'area se cancela y queda
\begin{equation}
  c\,\Delta t\;<\;\sqrt{L_p\,L_{3D}} ,
  \label{eq:estab}
\end{equation}
una condici\'on tipo CFL transparente: la onda de presi\'on no debe recorrer la
media geom\'etrica de la longitud del conducto y la del dominio en un paso de
tiempo. Con $c=8$~m/s, $\Delta t=10^{-3}$~s, $L_p=10$~mm y $L_{3D}\approx50$~mm
el margen es $\approx2.8$. Es m\'as ajustado que antes ---consecuencia de que
$C_a$ ya no es un par\'ametro libre--- y se reporta en la calibraci\'on.

\section{Calibraci\'on}
\label{sec:calib}

Se ejecuta una sola vez con la viscosidad newtoniana de referencia $\mu_N$ y se
congela, de manera que un cambio de reolog\'ia modifica el flujo y nunca la
geometr\'ia.

\paragraph{Lo que cambia.} Nada absorbe un residuo. Para cada \'arbol, las dos
relaciones de Poiseuille
\begin{equation}
  N\pi r^{2}\bar v = Q,
  \qquad
  \Delta P=\frac{8\mu_N L\,\bar v}{r^{2}}
  \label{eq:cierre}
\end{equation}
dejan tres grados de libertad entre $\{r,L,N,\bar v,\Delta P\}$. Se prescriben
tres observables fisiol\'ogicos ---velocidad media $\bar v$, tiempo de
tr\'ansito medio $T$, y la cuota medida $\Delta P$ del presupuesto de
presi\'on--- y el resto queda cerrado:
\begin{equation}
  \boxed{\;
  L=\bar v\,T,
  \qquad
  r=\bar v\sqrt{\frac{8\mu_N T}{\Delta P}},
  \qquad
  N=\frac{Q}{\bar v\,\pi r^{2}},
  \qquad
  V_0=N\pi r^{2}L=Q\,T\;}
  \label{eq:calib}
\end{equation}
La \'ultima igualdad es la relaci\'on de Stewart--Hamilton: el volumen del
compartimento es el volumen de indicador-diluci\'on, y los tiempos de
tr\'ansito coronarios son medibles directamente. El calibre $r$ deja de ser un
dato y pasa a ser una \emph{predicci\'on contrastable con morfometr\'ia}.

\paragraph{Pasos.}
\begin{enumerate}
\item \emph{Anatom\'ia y conducto.} $r_{p,i}=\sqrt{A_i/\pi}$ medido sobre la
      malla, $w_i=r_{p,i}^{3}/\sum_j r_{p,j}^{3}$, $Q_i=Q_{\text{tot}}w_i$,
      $\Delta P_{p,i}=8\mu_N L_{p,i}Q_i/(\pi r_{p,i}^{4})$; $Z_c$ y $C_a$ por
      \eqref{eq:zc}, $\Delta P_{z}=Z_cQ_i$.
\item \emph{Reparto medido del presupuesto.}
      $\Delta P_\mu=\Delta P-\Delta P_p-\Delta P_z$ con
      $\Delta P=p_{ar}(0)-P_{RA}$, y
      $\Delta P_v=f_v\Delta P_\mu$, $\Delta P_a=(1-f_v)\Delta P_\mu$.
\item \emph{\'Arboles.} \eqref{eq:calib} para cada uno.
\item \emph{\'Areas no tensionadas.} Con
      $\ptm^{\text{rep}}=P_{RA}+\Delta P_v-\alpha p_{LV}(0)$ se resuelve
      $a^{\text{rep}}$ en \eqref{eq:tubo} y se toma
      $A_u=\pi r^{2}/a^{\text{rep}}$, de modo que la ley de tubo reproduce
      exactamente el calibre calibrado en reposo. N\'otese $a^{\text{rep}}>1$:
      el lecho en reposo est\'a distendido por encima de su estado no
      tensionado, y eso es precisamente lo que le da algo que expulsar.
\end{enumerate}

\paragraph{Resultado.} Con $\bar v_a=5$~mm/s, $T_a=1.5$~s, $\bar v_v=3$~mm/s,
$T_v=2$~s y $f_v=0.13$:
\begin{center}\small
\begin{tabular}{lccl}
\toprule
Cantidad & Predicci\'on & Rango anat\'omico & Estado\\
\midrule
$r_a$ & $11.6~\mu$m & arteriola terminal, $8$--$15~\mu$m & compatible\\
$r_v$ & $20.7~\mu$m & v\'enula, $15$--$40~\mu$m & compatible\\
$L_a$ & $7.5$~mm & camino arteriolar & plausible (antes $46$~mm)\\
$L_v$ & $6.0$~mm & camino venular & plausible (antes $15$~mm)\\
$V_0$ & $7.0$~mL & volumen sangu\'ineo coronario LCA & compatible\\
$C$   & $3.9\times10^{-10}$~m$^3$/Pa & identificaciones $\sim10^{-10}$ & mismo orden\\
$\tau$ & $\approx0.4$~s & vaciado diast\'olico $0.2$--$0.5$~s & compatible\\
\bottomrule
\end{tabular}
\end{center}

\noindent
Que el calibre predicho caiga en la arteriola terminal y no en la arteriola de
$25~\mu$m que antes se \emph{imponia} es informativo: con $r_a=25~\mu$m fijo, las
relaciones \eqref{eq:cierre} dan
$\Delta P_a=8\mu\bar v^{2}T_a/r_a^{2}$, y exigir simult\'aneamente el $87\%$ del
presupuesto, $5$~mm/s y un tr\'ansito fisiol\'ogico es imposible: hab\'ia que
estirar $L_a$ hasta $46$~mm. La incompatibilidad estaba en los datos, no en el
esquema num\'erico.

\section{Constantes}

\begin{center}\small
\begin{tabular}{llll}
\toprule
S\'imbolo & C\'odigo & Valor & Origen\\
\midrule
$\mu_0,\mu_\infty$ & \texttt{viscosity.mu0,.muInf} & 0.353,\;4.181e-3 Pa\,s & reolog\'ia medida\\
$\lambda,n,a$ & \texttt{viscosity.lambda,.n,.yasuda} & 15.68 s,\;0.205,\;0.650 & reolog\'ia medida\\
$\mu_N$ & \texttt{newtonianCalibrationViscosity} & 3.5e-3 Pa\,s & referencia de calibraci\'on\\
\midrule
$r_p$ & \texttt{geometricParam.Rp} & de la malla & anatom\'ia (1.5--3.3 mm)\\
$L_p$ & \texttt{geometricParam.Lp} & 10--12.5 mm & anatom\'ia\\
$\bar v_a$ & \texttt{arteriolarVelocity} & 5 mm/s & fisiolog\'ia (2--10 mm/s)\\
$\bar v_v$ & \texttt{venousVelocity} & 3 mm/s & fisiolog\'ia (1--5 mm/s)\\
$T_a$ & \texttt{arteriolarTransitTime} & 1.5 s & tr\'ansito medido\\
$T_v$ & \texttt{venularTransitTime} & 2.0 s & tr\'ansito medido\\
$f_v$ & \texttt{venularPressureFraction} & 0.13 & reparto medido de carga\\
$K_a$ & \texttt{arteriolarWallStiffness} & 500 Pa & $E(h/r)^3$, morfometr\'ia\\
$K_v$ & \texttt{venularWallStiffness} & 200 Pa & $E(h/r)^3$, morfometr\'ia\\
$Q_{\text{tot}}$ & \texttt{lcaTargetFlow} & 2.0e-6 m$^3$/s & 120 mL/min\\
$c$ & \texttt{pulseWaveVelocity} & 8 m/s & onda de pulso (5--10)\\
$P_{RA}$ & \texttt{rightAtrialPressure} & 600 Pa & $\approx$4.5 mmHg\\
$\alpha$ & \texttt{intramyocardialFraction} & 0.65 & transmisi\'on tisular\\
\bottomrule
\end{tabular}
\end{center}

\noindent
Doce constantes libres al margen de la reolog\'ia y de $\mu_N$, frente a quince
antes. Desaparecen $r_a$, $r_v$, $L_v$, $C_{\text{tot}}$, $C_a/C$, $p_0$,
$\alpha_a$ y $\tau_f$; aparecen $T_a$, $T_v$, $f_v$, $K_a$ y $K_v$. Pero la
diferencia importante no es el recuento sino el cambio de estatus: $r_a$, $r_v$,
$L_a$, $L_v$, $N_a$, $N_v$, $V_0$, $C$, $C_a$ y $\tau$ pasan de ser entradas a
ser \emph{salidas}, y por tanto contrastables. Ninguna de las constantes
restantes es una resistencia ni una compliancia equivalente.

\section{Sensibilidad}
\label{sec:sens}

Barrido sobre el rango fisiol\'ogico de cada constante, con el modelo 0D
independiente conducido por las trazas $p_{ar}(t)$ y $p_{LV}(t)$ de una
simulaci\'on 3D previa:

\begin{center}\small
\begin{tabular}{lrrrr}
\toprule
Caso & $\bar Q$ (mL/min) & S/D & pico/media venoso & $V_{\min}$ (mL)\\
\midrule
referencia            & 66.0 & 0.54 & 4.27 & 4.49\\
$K_v=100$ / $400$~Pa  & 69.7 / 53.7 & 0.53 / 0.56 & 4.22 / 4.37 & 4.56 / 4.29\\
$K_a=250$ / $1000$~Pa & 47.6 / 76.8 & 0.56 / 0.54 & 4.08 / 4.35 & 4.12 / 4.73\\
$T_v=1.0$ / $3.0$~s   & 64.1 / 67.3 & 0.53 / 0.54 & 4.89 / 3.99 & 3.18 / 5.99\\
$f_v=0.10$ / $0.15$   & 62.8 / 67.7 & 0.54 / 0.53 & 4.51 / 4.14 & 4.17 / 4.68\\
$\alpha=0.50$ / $0.80$& 75.7 / 57.1 & \textbf{0.73 / 0.34} & 4.12 / 4.34 & 4.99 / 4.11\\
\bottomrule
\end{tabular}
\end{center}

\noindent
El cociente sist\'ole/di\'astole se mantiene en $0.53$--$0.56$ frente a
variaciones de factor cuatro en las rigideces de pared, de factor tres en el
tiempo de tr\'ansito venular y de $\pm20\%$ en el reparto de presi\'on. No es un
comportamiento de filo de navaja. La \'unica constante que lo mueve es $\alpha$,
y lo mueve de forma mon\'otona y fisiol\'ogicamente interpretable. $K_a$ es el
determinante del flujo medio, como corresponde al vaso de resistencia.

\section{Verificaci\'on}

El script \texttt{verify\_outlet\_0d.py} reimplementa la condici\'on de salida
por separado, la conduce con las trazas de una simulaci\'on previa sustituyendo
el dominio 3D por $p_{out}=p_{ar}(t)$, y ejecuta en paralelo la formulaci\'on
antigua y la nueva sobre la misma entrada:

\begin{center}\small
\begin{tabular}{lrr}
\toprule
 & Antigua & Nueva\\
\midrule
Jacobiano anal\'itico vs.\ diferencias finitas & $7\times10^{-8}$ & $4\times10^{-7}$\\
Iteraciones de Newton (media / m\'ax) & 3.2 / 7 & 3.6 / 6\\
Conservaci\'on discreta de volumen (relativa) & $10^{-13}$ & $4\times10^{-11}$\\
Volumen almacenado, m\'in--m\'ax (mL) & $0.00$--$0.41$ & $4.49$--$5.74$\\
Presi\'on transmural, m\'in--m\'ax (Pa) & $-8431$--$1281$ & $-378$--$-56$\\
\textbf{Sist\'ole/di\'astole del flujo de entrada} & \textbf{1.97} & \textbf{0.54}\\
Pico/media del flujo venoso & 5.74 & 4.27\\
\bottomrule
\end{tabular}
\end{center}

\noindent
El cociente sist\'ole/di\'astole entra en el rango fisiol\'ogico $0.4$--$0.6$;
era el defecto cualitativo pendiente. El lecho conserva una reserva de $4.5$~mL
en lugar de vaciarse por completo, y la presi\'on transmural se mantiene en una
banda estrecha en vez de recorrer $-8.4$~kPa.

\section{Lo que sigue sin cuadrar}

\paragraph{Pico venoso.} El cociente pico/media del flujo venoso baja de $5.74$
a $4.27$, pero el rango fisiol\'ogico es $2$--$3$. Es la \'unica m\'etrica que
sigue fuera de banda, y merece un tratamiento honesto: \emph{no responde a
ninguna constante del modelo}. Es insensible a $K_a$, $K_v$, $T_a$, $T_v$,
$f_v$, $\alpha$ y tambi\'en al caudal objetivo (probado hasta 216 mL/min, donde
sube a $4.40$ porque $V_0=QT$ escala con el caudal), y tampoco cede al suavizar
$p_{im}$. Un n\'umero que no responde a los par\'ametros no es un par\'ametro
mal puesto: es un elemento que falta. El candidato es la compliancia de las
venas extramioc\'ardicas y del seno coronario, que amortiguan la descarga
sist\'olica y que aqu\'i se sustituyen por una $P_{RA}$ r\'igida. A\~nadirla
ser\'ia el sim\'etrico exacto de $C_a$ en el lado arterial, y su valor saldr\'ia
de \eqref{eq:zc} aplicada al conducto venoso. Se deja fuera deliberadamente:
conviene comprobar primero que el diagn\'ostico es correcto.

\paragraph{Flujo medio y autorregulaci\'on.} $Q_{\text{tot}}$ se aplica en la
calibraci\'on como caudal del punto de reposo, y el modelo predice ahora un
d\'eficit sist\'olico, de modo que la media del ciclo cae a $\approx55\%$ del
objetivo. No es un error del esquema: es que un lecho coronario real
autorregula, ajustando el tono para mantener el flujo medio, y esto no se
modela. Caben dos lecturas coherentes: interpretar \texttt{lcaTargetFlow} como
el caudal diast\'olico de operaci\'on, o cerrar una iteraci\'on exterior sobre
el presupuesto para que la media del ciclo alcance el objetivo. La segunda
sigue siendo un punto fijo sobre una cantidad predicha, no un ajuste a la forma
de onda, y desplaza $r_a$ de $11.6$ a $12.3~\mu$m, es decir, nada.

\paragraph{Presi\'on motriz.} La onda diast\'olica depende tanto de $\tau$ como
de la evoluci\'on de $p_{ar}(t)$ y $p_{im}(t)$. Vale la pena se\~nalar que la
traza de $p_{LV}$ del modelo 0D alcanza $\dot p_{LV}=3.6\times10^{5}$~Pa/s,
frente a un valor fisiol\'ogico de $\sim2\times10^{5}$~Pa/s: el bloque
ventricular sube casi el doble de r\'apido de lo debido. Puesto que la amplitud
de la bomba ya no es un par\'ametro libre, cualquier exceso residual de
transitorios sist\'olicos debe buscarse ah\'i.

\paragraph{Alcance reol\'ogico.} Los efectos no newtonianos de la
microcirculaci\'on asociados a la naturaleza particulada de la sangre
(F\aa hr\ae us--Lindqvist) no est\'an contenidos en \eqref{eq:cy}. Con los
calibres ahora predichos ($11.6$ y $20.7~\mu$m) la omisi\'on es m\'as visible
que antes, porque son precisamente las escalas donde ese efecto opera.

\section{Criterio para distinguir correcci\'on de ajuste}

Se conserva el criterio de la formulaci\'on anterior y se aplica a lo nuevo.
Un t\'ermino leg\'itimo responde afirmativamente a: (i) \emph{¿su valor procede
de una medida independiente o de un primer principio?}; (ii) \emph{¿el
comportamiento es robusto en todo el rango fisiol\'ogico del par\'ametro?};
(iii) \emph{¿es defendible la ausencia del t\'ermino?}; (iv) \emph{¿se habr\'ia
incluido antes de observar el resultado?}

La ley de tubo \eqref{eq:tubo} las satisface: $K=E(h/r)^3/12(1-\nu^2)$ es
material y morfom\'etrico; el barrido de \S\ref{sec:sens} muestra robustez de
factor cuatro en $K$; un vaso de pared delgada rodeado de m\'usculo en
contracci\'on cuyo calibre no dependa de la presi\'on transmural es
f\'isicamente insostenible; y deber\'ia haberse escrito as\'i desde el
principio, puesto que ya se admit\'ia que el volumen depend\'ia de $\ptm$.

La derivada \eqref{eq:dqdr} no es un t\'ermino de modelo sino una identidad, y
por tanto no se somete al criterio; pero conviene notar que sin ella la
alternativa habr\'ia sido un Jacobiano aproximado, y con \'el la tentaci\'on de
retocar la amortiguaci\'on hasta que convergiera.

En cambio, la eliminaci\'on de $\tau_f$ y $\alpha_a$ se justifica por el mismo
criterio le\'ido al rev\'es: no ten\'ian medida independiente, se eleg\'ian
mirando el resultado, y su ausencia es ahora perfectamente defendible porque el
mecanismo que ven\'ian a paliar ya no existe.

\end{document}
